\section{Bachelor Angewandte Informatik}
\sidebar{
    \centering
    \includegraphics[width=3cm]{bilder/cant_sleep_1.png}\\\vspace{13mm}
    \includegraphics[width=3cm]{bilder/cant_sleep_2.png}\\\vspace{13mm}
    \includegraphics[width=3cm]{bilder/cant_sleep_3.png}\\\vspace{13mm}
    \includegraphics[width=3cm]{bilder/cant_sleep_4.png}
}

Im ersten Semester hört ihr die Informatikvorlesungen „Einführung in die Praktische Informatik“
und „Einführung in die Technische Informatik“ sowie eine Mathematikvorlesung.
Auch ist ein Programmierkurs Pflicht, in dem grundlegende Fertigkeiten in C++ vermittelt werden
sollen. Außerdem nehmt ihr an einem, von der Fakultät organsierten Mentoringprogramm, teil -- das Modul
heißt „Einführung in das Studium“ --, das euch den Einstieg ins Studium erleichtern soll.

In „Einführung in die Praktische Informatik“ (auch „Informatik~I“ genannt) sollt ihr Programmieren
und andere Grundfertigkeiten lernen. Ihr sollt einen Überblick über verschiedene Grundkonzepte der
Informatik und Grundkenntnisse in einer oder mehrerer Programmiersprachen bekommen.

Die „Einführung in die Technische Informatik“ („Informatik~II“) soll euch Kenntnisse über den
grundsätzlichen Aufbau und Funktionsweise von Rechnern vermitteln. Ihr lernt Schritt für Schritt,
angefangen bei Logik und einfachen Schaltungen, wie ein Prozessor funktioniert.

In den späteren Vorlesungen werden einzelne Themengebiete eröffnet und vertieft, die Namen der
Vorlesungen sprechen zum größten Teil für sich, zum Beispiel „Datenbanken“ oder „Betriebssysteme
und Netzwerke“.

Das Informatikstudium in Heidelberg beinhaltet einen in Vergleich zu anderen Universitäten hohen
Anteil an Mathematik. Ihr könnt von Anfang an entscheiden, welche Art der Mathematikausbildung ihr
erhalten möchtet. Ihr habt die Wahl zwischen mehreren Varianten.

In der ersten Variante hört ihr die Vorlesungen Lineare Algebra I (im 1.
Semester) und Analysis I (im 3. Semester). Alternativ könnt ihr als Ersatz für
die „Lineare Algebra I“ und die „Analysis I“ im ersten und zweiten Semester die
Veranstaltungen „Mathematik für Informatiker I und II“ hören.  Diese
Vorlesungen soll den Informatikstudierenden speziell die Mathematikkenntnisse
vermitteln, die sie in ihrem späteren Studium brauchen werden. Diese
Veranstaltung ist relativ neu, sie bietet aber trotzdem solide
Mathematikgrundlagen und den Bezug zur Informatik.

Hinzu kommen noch die Vorlesung Einführung in die Numerik und mindestens ein
Mathematikmodul eurer Wahl, wobei mindestens Analysis 2, Mathematische Logik
oder Einführung in die Wahrscheinlichkeitstheorie und Statistik belegt werden
muss.  Maximal dürft ihr aber insgesamt 5 Mathematikvorlesungen belegen, d.\
h.\ zu den oben beschriebenen Modulen dürft ihr maximal ein weiteres
Mathematikmodul belegen.

Die Mathematikmodule unterscheiden sich stark von der Schulmathematik,
unterschätzt den Aufwand für die Vorlesungen nicht!  Wenn ihr schon wisst, dass
ihr euer Informatikstudium mathematisch ausrichten und weitere Mathematikmodule
hört, euch vielleicht später sogar im wissenschaftlichen Rechnen spezialisieren
möchtet, ist die Variante mit Linarer Algebra I und Analysis I empfehlenswert.

Zum Thema Prüfungen und Prüfungswiederholungen: Um in den Vorlesungen
Creditpoints (meisten mit CP oder LP abgekürzt) zu bekommen, müsst ihr am Ende
in den meisten Fällen eine schriftliche Prüfung bestehen. Die genauen
Modalitäten, um überhaupt für die Prüfung zugelassen zu werden, legt die
Dozentin bzw.\ der Dozent am Anfang fest.  Informiert euch unbedingt genau,
worin die bestehen! Meist müsst ihr am Ende nur im Mittel 50 oder 60 Prozent
der Punkte auf den Übungszetteln erreichen, manchmal ist aber auch gefordert,
dass jeder Übungszettel bearbeitet wurde.

Ihr solltet euch auch bei jeder Veranstaltung genau darüber informieren, was
genau \emph{eine} Prüfung beinhaltet (z.\,B.\ Bestehen von einer von zwei
Klausuren). Insbesondere der letzte Teil ist wichtig, denn ihr könnt Prüfungen
grundsätzlich zweimal versuchen. Je nach Veranstaltung und Dozent/in
\emph{können} zwei Klausuren als eine Prüfung zählen, müssen aber nicht -- dann
würde jede geschriebene Klausur als ein Prüfungsversuch gelten. Wenn ihr aus
Gründen, die ihr nicht zu verantworten habt (wie krank sein), nicht an einer
\emph{Prüfung} teilnehmen konntet, erhöhen sich entsprechend eure Versuche.
Aber keine Panik -- auch wenn ihr in einer Vorlesung zum zweiten Mal die
Prüfung nicht bestanden haben solltest, seid ihr noch nicht exmatrikuliert: In
bis zu vier Fällen dürft ihr auf Antrag Prüfungen auch ein zweites Mal
wiederholen, dies gilt aber nicht für die Orientierungsprüfung und für die
Bachelorarbeit.

Denkt dran: Eine Klausuranmeldung ist eine Anmeldung zu einer Prüfung. Wenn ihr
euch sicher seid, sie nicht zu bestehen, überlegt euch lieber zweimal, ob ihr
euch anmeldet. Das Schöne an eurem Studiengang ist, dass die Noten der
Grundpflichtmodule nicht in eure Abschlussnote zählen, das heißt, dass auch
wenn die ersten Semester mit der ganzen Mathematik vielleicht etwas hart sind,
eure Abschlussnote darunter nicht unbedingt zu leiden hat.

Auch in Informatik gibt es eine sog. Orientierungsprüfung. Die besteht darin,
dass ihr den Schein der Vorlesung „Einführung in die Praktische Informatik“
bekommt. Diesen Schein müsst ihr spätestens am Ende des 3. Semester vorweisen
können.

Programmieren ist eine wichtige Fertigkeit. Und die werdet ihr nur teilweise an
der Uni lernen. Im Programmierkurs und in der Info I wird versucht, euch die
Konzepte näher zu bringen und in den Übungsaufgaben werdet ihr selber kleinere
Sachen programmieren müssen. Aber so lernt ihr letztlich nicht, wie man
„richtig“ programmiert.

Es gibt eine enorme Anzahl an Programmiersprachen und die meisten haben ihre
Daseinsberechtigung: um praktische Probleme zu lösen, akademisches Interesse
oder auch Unterhaltung. Es gibt diverse Paradigmen, nach denen
Programmiersprachen konzipiert sind und an der Uni werdet ihr nur die wenigsten
kennen lernen. Auch solltet ihr das System kennen, auf dem ihr arbeitet: solide
Kenntnisse von unixoiden Systemen sind immer Gold wert.

Die Uni hilft aber durchaus: Die „Einführung in die Praktische Informatik“
beinhaltet eine Art Programmierkurs, es gibt einen richtigen Programmierkurs,
der auch Pflicht ist, zur „Numerik 0“ gibt es praktische Übungen, in denen
ebenfalls programmiert wird. Eigeninitiative ist aber dennoch wichtig.

\vspace{-\parskip}
